\documentclass[12pt, a4paper]{article}

% --- PACKAGES ---
\usepackage[utf8]{inputenc}
\usepackage[margin=1in]{geometry}
\usepackage{setspace}
\usepackage{graphicx}
\usepackage{natbib} % <--- Added this to fix the square brackets!
\usepackage{hyperref}
\usepackage{titlesec}
\usepackage{booktabs}
\usepackage{tabularx}

% --- FORMATTING ---
\onehalfspacing
\hypersetup{
    colorlinks=true,
    linkcolor=black,
    filecolor=black,      
    urlcolor=blue,
    citecolor=black, % <--- Added this to kill the neon green!
    pdftitle={Dual-Use ICT Facilities -- Maigo},
}

\begin{document}

% --- TITLE BLOCK ---
\begin{center}
    \Large\textbf{Institutionalizing Dual-Use ICT Facilities: A Public-Private Partnership for Technical-Vocational Education in the Municipality of Maigo} \\[1.5em]
    \large \textbf{Jemar John J Lumingkit\textsuperscript{1}} \\[0.5em]
    \normalsize \textsuperscript{1}Mindanao State University – Iligan Institute of Technology \\
    *Iligan City, Philippines
\end{center}
\vspace{1em}

% --- ABSTRACT ---
\begin{abstract}
\noindent This policy proposal addresses a critical infrastructure paradox in Maigo, Lanao del Norte: a severe deficit of Information and Communication Technology (ICT) laboratories for Technical-Vocational Education and Training (TVET) programs at MSU-MCEST and TESDA, juxtaposed against the daytime underutilization of local internet cafés. While public institutions face prohibitive capital expenditure (CapEx) constraints, local Micro, Small, and Medium Enterprises (MSMEs) possess high-performance computing assets that sit idle during academic hours.

Anchored in Collaborative Governance and the Philippine Public-Private Partnership (PPP) Code (RA 11966), this study proposes the institutionalization of ``Dual-Use ICT Facilities.'' This service-contracting framework enables the Local Government Unit (LGU) to lease private cybercafés as accredited training micro-laboratories. To ensure pedagogical integrity and cybersecurity, the proposal mandates a Diskless (PXE-Boot) Multiple Image architecture. This mechanism deploys a restricted ``Educational Image'' during school hours---effectively segregating recreational software---before reverting to ``Commercial Mode'' for nighttime operations.

Employing a qualitative case study design, this research evaluates the legal, economic, and technical feasibility of the proposed framework. The ultimate output is a comprehensive Municipal Ordinance that delivers a Zero-CapEx solution for public education, stabilizes MSME revenue streams, and establishes a scalable blueprint for rural digital resilience.

\vspace{0.5em}
\noindent \textbf{Keywords---}Dual-Use, Public-Private Partnership (PPP), Technical-Vocational Education (TVET), MSME Resilience, Collaborative Governance, Diskless Architecture.
\end{abstract}

\newpage

% --- MODULAR CHAPTER INCLUSIONS ---
\section*{Executive Summary}
\addcontentsline{toc}{section}{Executive Summary}

\textbf{To:} The Office of the Mayor and the Sangguniang Bayan of Maigo \\
\textbf{Subject:} Institutionalizing Dual-Use ICT Facilities via Public-Private Partnership \\
\textbf{Framework:} RA 11966 (PPP Code of the Philippines)

\vspace{1em}
\noindent \textbf{The Core Problem: Infrastructure Scarcity vs. Asset Underutilization} \\
The Municipality of Maigo faces a critical bottleneck in its digital upskilling initiatives. Key educational institutions, including MSU-MCEST and local TESDA branches, lack the capital expenditure (CapEx) required to build and maintain high-end computer laboratories for Technical-Vocational Education and Training (TVET). Consequently, residents are deprived of advanced ICT certifications. Simultaneously, local Micro, Small, and Medium Enterprises (MSMEs) operating internet cafés possess high-performance hardware that sits idle and depreciates during daytime academic hours due to shifts in mobile internet usage.

\vspace{1em}
\noindent \textbf{The Proposed Solution: The "Dual-Use" PPP Model} \\
This policy research proposes the institutionalization of a service-contracting model under the 2023 PPP Code. Instead of purchasing redundant hardware, the Local Government Unit (LGU) will lease accredited private cybercafés during standard academic hours (8:00 AM -- 5:00 PM) to serve as public training micro-laboratories. 

To ensure a secure and distraction-free learning environment, the policy mandates a \textbf{Diskless (PXE-Boot) Multiple Image architecture}. This technical requirement allows the hardware to boot into a sterile, government-controlled ``Educational Image'' (containing only TVET software and blocking all games/social media) during the day, before reverting to standard commercial operations at night.

\vspace{1em}
\noindent \textbf{Strategic Benefits to the Municipality of Maigo:}
\begin{itemize}
    \item \textbf{Zero-CapEx Upgrades:} The LGU and public institutions bypass the multi-million peso costs of procuring hardware, facility construction, and IT maintenance.
    \item \textbf{MSME Economic Resilience:} Struggling local internet cafés receive a stable, institutionalized revenue stream through daytime government leasing, preventing business closures.
    \item \textbf{Bridging the Digital Divide:} Out-of-school youth and students gain immediate access to industry-standard hardware for advanced certifications (e.g., AutoCAD, CSS NC II, Programming).
\end{itemize}

\vspace{1em}
\noindent \textbf{Actionable Output} \\
This document outlines the comparative policy alternatives and provides a draft \textbf{Municipal Ordinance} and \textbf{Service Level Agreement (SLA)} to legally and technically operationalize this framework.
\newpage
\section*{CHAPTER 1: INTRODUCTION}
\addcontentsline{toc}{section}{CHAPTER 1: INTRODUCTION}

\subsection*{1.1 Background of the Study}
The rapid acceleration of digital technology has created a socio-economic paradox in developing municipalities. While digital literacy is now a mandatory prerequisite for global and local employment, the infrastructure required to teach these skills is lagging. In many areas, the Technical Education and Skills Development Authority (TESDA) and local universities are under immense pressure to provide computer-based certifications, yet they are hamstrung by acute capital scarcity. Conversely, the widespread accessibility of smartphones has triggered an economic decline in the traditional internet café (PC bang) industry, leaving high-specification hardware idle during the day.

In the Municipality of Maigo, Lanao del Norte, this paradox is highly pronounced. The municipality hosts the Mindanao State University - Maigo College of Education, Science and Technology (MSU-MCEST) and local TESDA operations, both of which face significant infrastructural constraints. MSU-MCEST’s computer laboratories primarily cater to basic introductory courses, lacking the capacity for advanced, certification-level training. Similarly, local TESDA programs cannot offer crucial ICT certifications because they lack the specialized laboratories required for training and assessment.

Simultaneously, Maigo’s internet café owners are facing bankruptcy. As their traditional demographic—students seeking basic research or entertainment—migrates to mobile data, their high-performance hardware sits dormant during standard business hours (8:00 AM to 5:00 PM), accruing depreciation without generating revenue.

This study posits that these isolated struggles represent a significant opportunity for Collaborative Governance. By institutionalizing a Public-Private Partnership (PPP) that transforms struggling internet cafés into government-contracted, ``Dual-Use Tech Spaces,'' the LGU can bridge the digital divide through a ``Zero-CapEx'' model while providing an economic lifeline to local MSMEs.

\subsection*{1.2 Statement of the Problem}
This study aims to develop a localized public policy framework to institutionalize the dual-use of commercial internet cafés as government-contracted training micro-laboratories during the day and commercial hubs at night in the Municipality of Maigo.

Specifically, it seeks to answer the following:
\begin{enumerate}
    \item \textbf{The Governance and Legal Factor:} What existing local ordinances, zoning laws, or procurement guidelines in Maigo currently hinder or support the LGU in formally leasing private IT infrastructure?
    \item \textbf{The Business and Economic Factor:} What financial models (e.g., service contracting, utility subsidies) are required to make government leasing a viable economic lifeline for Maigo’s MSMEs?
    \item \textbf{The Infrastructure and Technical Factor:} How do the existing hardware and network architectures (such as Diskless/PXE-Boot systems) in local cafés measure against the baseline standards mandated by TESDA and MSU-MCEST?
    \item \textbf{The Policy Output:} Based on the findings, what comprehensive LGU ordinance and implementation plan can be proposed to operationalize the ``Dual-Use Tech Space'' program?
\end{enumerate}

\subsection*{1.3 Objectives of the Study}
\textbf{General Objective:} To design a formalized PPP framework and municipal ordinance that utilizes private internet cafés in Maigo as official TESDA and university-accredited computer laboratories.

\textbf{Specific Objectives:}
\begin{itemize}
    \item To assess the current IT infrastructure gap within MSU-MCEST and TESDA centers in Maigo.
    \item To evaluate the technical readiness and economic viability of local internet cafés to serve as formal educational spaces.
    \item To identify the legal and bureaucratic mechanisms under the Philippine PPP Code (RA 11966) for daytime leasing agreements.
    \item To draft a localized policy recommendation that addresses the digital divide through collaborative asset utilization.
\end{itemize}

\subsection*{1.4 Significance of the Study}
\begin{itemize}
    \item \textbf{The Local Government of Maigo:} Provides a Zero-Capital-Expenditure (Zero-CapEx) model to upgrade local digital infrastructure without building redundant facilities.
    \item \textbf{Internet Café Owners (MSMEs):} Offers a sustainable, institutionalized revenue stream to maximize hardware ROI and prevent business closure.
    \item \textbf{TESDA and MSU-MCEST:} Grants immediate access to high-spec, maintained laboratories, expanding their capacity to offer advanced ICT certifications.
    \item \textbf{The Citizens of Maigo:} Democratizes access to digital skills training, directly combating the digital divide and improving employability for out-of-school youth and workers.
\end{itemize}

\subsection*{1.5 Scope and Delimitations}
This study focuses strictly on the public administration and policy formulation aspects of utilizing private internet cafés for public education in Maigo. The subjects include registered café owners, representatives from MSU-MCEST, TESDA administrators, and LGU officials. The research will not assess pedagogical effectiveness or conduct the actual training; it focuses exclusively on the governance mechanism of providing the space.

\subsection*{1.6 Conceptual Framework}
The study is anchored on Collaborative Governance and New Public Management (NPM). The LGU acts as the ``Central Orchestrator,'' using its legislative power to create accreditation systems and its executive power to facilitate service contracts. This bridges the ``Educational Deficit'' (lack of hardware) with the ``Private Sector Surplus'' (underutilized hardware), resulting in an institutionalized policy that serves the public good while sustaining private enterprise.
\newpage
\section*{CHAPTER 2: REVIEW OF RELATED LITERATURE AND STUDIES}
\addcontentsline{toc}{section}{CHAPTER 2: REVIEW OF RELATED LITERATURE AND STUDIES}

\subsection*{2.1 The Digital Divide and Infrastructure Deficit}
The ``digital divide'' in the Philippines is often a matter of hardware access rather than just connectivity. Technical learning (e.g., programming, CAD, graphics) requires desktop computers, which public institutions struggle to procure and maintain due to rapid depreciation and high CapEx requirements \cite{ulas2018}.

\subsection*{2.2 The Economic Decline of the Internet Café Industry}
Market analyses indicate a severe decline in the internet café sector as users shift to mobile platforms \cite{garcia2019}. While these shops possess high-end hardware for esports, they face heavy downtime during the day, making them prime candidates for government resource-sharing models.

\subsection*{2.3 Cybercafes as Informal Learning Spaces (ICT4D)}
Literature in ICT for Development (ICT4D) recognizes cybercafes as vital community hubs \cite{robinson2009}. However, they are currently used as informal student workarounds. This study proposes moving from informal usage to institutionalized governance.

\subsection*{2.4 The Philippine PPP Code (RA 11966)}
The newly enacted PPP Code of 2023 \cite{pppcode2023} recognizes the autonomy of LGUs to enter into local PPP projects. It emphasizes Value for Money (VFM), allowing LGUs to lease private assets to deliver public services (digital education) more efficiently than traditional procurement.

\subsection*{2.5 TESDA Modalities and Precedents}
TESDA’s Community-Based Training (CBT) and Educational Service Contracting (ESC) provide precedents for the government paying private partners to deliver public education. This study applies this ``infrastructure service contracting'' logic to the ICT sector.
\newpage
\section{Methodology}

\subsection*{3.1 Research Design}
This study employs a Qualitative Single Instrumental Case Study design. This approach is ideal for understanding the complex legal, technical, and administrative intersections required to draft a new municipal policy.

\subsection*{3.2 Research Locale}
The study is situated in the Municipality of Maigo, Lanao del Norte, a locale characterized by a strong presence of technical-vocational institutions (MSU-MCEST and TESDA) and a cluster of struggling commercial internet cafés.

\subsection*{3.3 Key Informants}
The researcher will utilize Purposive Sampling across three sectors:
\begin{itemize}
    \item \textbf{Governance:} LGU Officials (Mayor, SB Committee on Education, PESO Manager).
    \item \textbf{Educational:} Administrators from MSU-MCEST and TESDA Provincial/Municipal coordinators.
    \item \textbf{Private:} 3–5 registered internet café owners in Maigo.
\end{itemize}

\subsection*{3.4 Research Instruments}
Data will be collected via Semi-Structured Interview Guides and Documentary Analysis of local ordinances and TESDA facility standards.

\subsection*{3.5 Data Gathering and Analysis}
Data will be gathered through clearances, document retrieval, and interviews. The study will use Thematic Analysis and Triangulation to cross-reference business needs, legal boundaries, and educational standards to draft the final policy output.

\subsection*{3.6 Technical System Architecture (The ``Dual-Use'' Mechanism)}
For this policy to be accepted, the researcher must demonstrate how the commercial and educational environments are separated. The study proposes the utilization of Diskless (PXE Boot) Architecture, which is already common in high-end Philippine internet cafes.
\begin{itemize}
    \item \textbf{Environmental Segregation:} The study will explore how a single server can host two distinct ``Images'':
    \begin{itemize}
        \item \textit{The ``Edu-Image'' (Daytime):} A restricted OS containing only Microsoft Office, AutoCAD, Programming IDEs, and TESDA assessment tools. All gaming ports and social media are blocked at the BIOS/Network level.
        \item \textit{The ``Gamer-Image'' (Nighttime):} The standard commercial OS with high-fidelity games and unrestricted internet access.
    \end{itemize}
    \item \textbf{Data Volatility:} The research will investigate the ``Reboot-to-Restore'' feature, ensuring that any files or potential malware introduced during one session are completely wiped before the next session begins.
\end{itemize}

\subsection*{3.7 Proposed Implementation Roadmap (The Policy Lifecycle)}
The study will outline a four-phase transition for the Municipality of Maigo to adopt this framework:

\begin{table}[h!]
\centering
\begin{tabular}{|l|p{7cm}|p{5cm}|}
\hline
\textbf{Phase} & \textbf{Activity} & \textbf{Responsible Entities} \\ \hline
I. Vetting & Technical and physical audit of local cafes based on TESDA standards. & LGU, TESDA, MSU-MCEST \\ \hline
II. Accreditation & Granting of ``Certified Dual-Use Tech Space'' status to eligible MSMEs. & Maigo Municipal Council \\ \hline
III. Integration & Deployment of the ``Edu-Image'' software and instructor training. & University IT Dept / TESDA \\ \hline
IV. Operation & Commencement of daytime classes and monthly lease disbursement. & LGU Finance / PESO \\ \hline
\end{tabular}
\end{table}

\subsection*{3.10 Anticipated Policy Outcomes}
\begin{itemize}
    \item \textbf{Economic Resilience:} A 30–40\% increase in daytime revenue for participating MSMEs, preventing local business closures.
    \item \textbf{Educational Expansion:} A 50\% increase in the capacity of MSU-MCEST and TESDA to accept ICT scholars without new building construction.
    \item \textbf{Governance Innovation:} A scalable ``Maigo Model'' of Public-Private Partnership that can be replicated by other 4th and 5th-class municipalities in Lanao del Norte.
\end{itemize}
\newpage
\section{Results and Discussion}
To determine the most viable approach for addressing the ICT infrastructure deficit in the Municipality of Maigo, this study evaluates three distinct policy alternatives. These alternatives are assessed based on three criteria: Economic Feasibility (Capital Expenditure vs. Operating Expenditure), Technical Feasibility (Maintenance and longevity), and Socio-Economic Impact (Value for Money and MSME resilience).

\subsection{Alternative A: Do nothing (Status Quo)}
Under this alternative, the Local Government Unit (LGU) and educational institutions maintain their current operations. MSU-MCEST and TESDA continue to rely on existing, limited laboratory space, while out-of-school youth remain without access to advanced ICT certifications.

\begin{itemize}
    \item Economic Impact: While this requires zero government spending initially, the opportunity cost is severe. The municipality loses potential skilled labor.
    \item Socio-Economic Impact: Local MSMEs (internet cafés) will continue to face declining revenues and potential bankruptcy due to underutilized hardware during standard business hours.
    \item Verdict: Highly unsustainable and fails to address the digital divide.
\end{itemize}

\subsection{Alternative B: LGU builds new public labs (High CapEx)}
This is the standard approach where the LGU or national government funds the construction of new, dedicated public computer laboratories.

\begin{itemize}
    \item Economic Impact: This requires massive Capital Expenditure (CapEx). A standard 30-seat laboratory requires funding for high-end PCs, air conditioning, structural security, and commercial-grade networking equipment. Furthermore, IT hardware depreciates rapidly, requiring hardware replacement cycles every 3 to 5 years.
    \item Technical Impact: Requires the LGU or schools to hire dedicated, full-time IT administrators for continuous maintenance.
    \item Verdict: While it solves the educational deficit, it is economically prohibitive for a standard municipality and does nothing to support the local MSME sector.
\end{itemize}

\subsection{Alternative C: The Dual-Use PPP Model (Zero-CapEx)}
This alternative represents the proposed framework: institutionalizing a service-contracting model under the PPP Code (RA 11966) where the LGU leases existing private cybercafés during the day.

\begin{itemize}
    \item Economic Impact: Operates on a Zero-CapEx model. The government only incurs Operating Expenditures (OpEx) through hourly or monthly leasing. The financial risk of hardware depreciation and maintenance remains with the private sector.
    \item Socio-Economic Impact: Provides a guaranteed daytime revenue stream for struggling MSMEs, preventing business closures and stabilizing the local economy.
    \item Technical Impact: Utilizing Diskless (PXE-Boot) Multiple Image architecture ensures a sterile, distraction-free educational environment during the day without permanently altering the café's commercial capabilities at night.
    \item Verdict: Highly feasible. It leverages existing, idle community assets to deliver public services efficiently.
\end{itemize}

\subsection{4.4 Comparative Analysis Matrix}
The comparison heavily favors Alternative C as the most viable public policy intervention. By shifting from a "procurement of goods" mindset to a "procurement of services" mindset, the LGU achieves maximum Value for Money (VFM). Alternative C is the only option that simultaneously solves the educational infrastructure deficit while acting as an economic stimulus for local businesses.
\newpage
\section{The Policy Recommendations and Implementation}

\subsection{The Proposed Municipal Ordinance}
To formalize the Public-Private Partnership (PPP) and ensure continuous funding, the Sangguniang Bayan of Maigo must enact a local law. Below is the proposed legislative draft:

\vspace{1em}
\begin{center}
    \textbf{DRAFT MUNICIPAL ORDINANCE NO. \_\_\_ - 2026} \\
    \vspace{0.5em}
    \textbf{AN ORDINANCE INSTITUTIONALIZING THE DUAL-USE ICT FACILITY PROGRAM, ESTABLISHING A PUBLIC-PRIVATE PARTNERSHIP WITH LOCAL INTERNET CAFÉS FOR TECHNICAL-VOCATIONAL EDUCATION, APPROPRIATING FUNDS THEREFOR, AND FOR OTHER PURPOSES.}
\end{center}

\vspace{0.5em}
\noindent \textbf{WHEREAS,} Section 16 of Republic Act No. 7160 (Local Government Code) mandates local government units to promote the general welfare, enhance economic prosperity, and improve public morals;

\noindent \textbf{WHEREAS,} Republic Act No. 11966 (PPP Code of the Philippines) empowers LGUs to engage in service-contracting with the private sector to deliver public services efficiently;

\noindent \textbf{WHEREAS,} there is an acute shortage of capital-intensive ICT laboratories for TVET and university students in the Municipality of Maigo, alongside the severe daytime economic underutilization of local internet cafés (MSMEs);

\noindent \textbf{BE IT ORDAINED BY THE SANGGUNIANG BAYAN OF MAIGO:}

\vspace{0.5em}
\noindent \textbf{Section 1. Title.} This Ordinance shall be known as the \textit{"Maigo Dual-Use ICT Facility Act of 2026."}

\noindent \textbf{Section 2. Declaration of Policy.} It is the policy of the Municipality of Maigo to bridge the digital divide and support local MSME resilience through collaborative governance, transitioning from a procurement-heavy infrastructure model to a sustainable service-contracting model.

\noindent \textbf{Section 3. The Dual-Use Program.} The LGU shall formally lease accredited private internet cafés during standard academic hours (8:00 AM to 5:00 PM) to serve as public training micro-laboratories for MSU-MCEST and TESDA scholars. 

\noindent \textbf{Section 4. Accreditation and Technical Standards.} To qualify for the program, private MSMEs must utilize a \textbf{Diskless (PXE-Boot) Multiple Image architecture}. Participating facilities must maintain an isolated "Educational Image" that completely restricts access to recreational games and non-educational websites during leased hours.

\noindent \textbf{Section 5. Exemption from Anti-Truancy Apprehension.} Accredited facilities operating under the "Educational Image" during official leased hours are hereby exempted from daytime anti-truancy ordinances, provided that students present their official LGU or University access passes.

\noindent \textbf{Section 6. Appropriations.} The Municipal Government shall allocate an initial funding of \textit{[Insert Amount]} from the Special Education Fund (SEF) or the General Fund to cover the hourly lease vouchers for the first fiscal year of implementation.

\subsection{The Service Level Agreement (SLA) Guidelines}
The ordinance establishes the law, but the SLA establishes the daily business rules between the Maigo LGU and the internet café owner. The SLA must explicitly define the following operational boundaries:

\begin{itemize}
    \item \textbf{LGU / Educational Institution Duties:} The school's IT Department is strictly responsible for building, updating, and patching the "Educational VHD (Virtual Hard Disk) Image." The LGU commits to paying lease vouchers within 15 working days of the billing cycle.
    \item \textbf{MSME (Café Owner) Duties:} The private partner is responsible for physical security, electricity, internet bandwidth, and ensuring the server boots the correct "Educational Image" before 8:00 AM. 
    \item \textbf{Purge Protocol:} The SLA mandates a daily "Reboot-to-Restore" or write-back purge. This ensures any user data, browsing history, or malware introduced by a student is completely wiped the moment the PC restarts, protecting both the student's privacy and the MSME's hardware.
\end{itemize}

\subsection{The Risk Allocation Matrix}
A core tenet of the Philippine PPP Code is equitable risk distribution. Table 1 outlines the allocation of responsibilities under the Dual-Use framework, proving to the Commission on Audit (COA) that the government is protected from private hardware liabilities.

\begin{table}[htbp]
\centering
\caption{Risk Allocation Matrix for the Dual-Use ICT Facility Framework}
\vspace{0.5em}
\begin{tabularx}{\textwidth}{|>{\raggedright\arraybackslash}p{2.5cm}|>{\raggedright\arraybackslash}p{4.5cm}|>{\raggedright\arraybackslash}p{2cm}|X|}
\hline
\textbf{Risk Category} & \textbf{Potential Event} & \textbf{Risk Owner} & \textbf{Mitigation \& SLA Protocol} \\
\hline
\textbf{Hardware \& Utility} & Server failure, internet outage, or power loss during class. & Private MSME & MSME must provide Uninterruptible Power Supplies (UPS). LGU receives pro-rated refunds for downtime. \\
\hline
\textbf{Software \& Cyber} & Corruption of the ``Educational Image'' or malware infection. & LGU / School IT & School IT maintains exclusive admin rights to the Educational VHD. Reboot-to-Restore purges daily risks. \\
\hline
\textbf{Operational} & Student damages a monitor, peripheral, or physical property. & LGU / School & LGU compensates the MSME based on a pre-agreed depreciation table, followed by school disciplinary action. \\
\hline
\textbf{Financial} & Delayed LGU lease payments to the MSME. & LGU & Standard PPP Code late payment interest applies to protect the small business owner. \\
\hline
\end{tabularx}
\end{table}
\newpage
\section{Conclusion}
Final thoughts and the step-by-step rollout plan for the Maigo LGU.
\newpage

% --- BIBLIOGRAPHY ---
\newpage
\addcontentsline{toc}{section}{REFERENCES}
\bibliographystyle{apalike}
\bibliography{bibliography_policy}

\end{document}